\documentclass[UTF8]{ctexart}
\newtheorem{thm}{定理}
\usepackage{graphicx}
\usepackage{geometry}
\usepackage{amsmath}
\usepackage{float}
\usepackage[format=hang,font=small,textfont=it]{caption}
\usepackage[nottoc]{tocbibind}
\title{\heiti 杂谈勾股定理}
\author{\kaishu 孙凯威}
\date{\today}
\bibliographystyle{plain}
\newenvironment{myquote}{\begin{quote}\kaishu\zihao{-5}}
{\end{quote}}
\geometry{a6paper,centering,scale=0.8}
\begin{document}
    \maketitle
    \begin{abstract}
        这是一篇关于勾股定理的小短文。
    \end{abstract}
    \tableofcontents

    \section{勾股定理在古代}
    西方称勾股定理为毕达哥拉斯定理，将勾股定
理的发现归功于公元前 6 世纪的毕达哥拉斯学派
\cite{Kline}。该学派得到了一个法则，可以求出可排成直角
三角形三边的三元数组。毕达哥拉斯学派没有书面
	著作，该定理的严格表述和证明则见于欧几里德\footnote{欧几里得，约公元前 330--275 年。}
《几何原本》的命题 47：“直角三角形斜边上的正方
形等于两直角边上的两个正方形之和。”证明是用面
积做的。

我国《周髀算经》载商高（约公元前 12 世纪）
答周公问：
    \begin{myquote}
	    勾广三，股修四，径隅五
    \end{myquote}
又载陈子（约公元前 7--6 世纪）答荣方问：
\begin{myquote}
	若求邪至日者，以日下为勾，日高为股，勾股各自乘，并而开方除之，得邪至日。
\end{myquote}
都较古希腊更早。后者已经明确道出勾股定理的一
般形式。图\ref{fig:xiantu}是我国古代对勾股定理的一种证明\cite{quanjing}。
\begin{figure}[ht]
    \centering
    \includegraphics[width=3cm]{xiantu.png}
    \caption{宋赵爽在《周髀算经》注中作的弦图（仿制），该图给出了勾股定
    5 理的一个极具对称美的证明。 }
    \label{fig:xiantu}
\end{figure}
\section{勾股定理的现在形式}
勾股定理可以用现代语言表述如下：
    \begin{thm}[勾股定理]
        直角三角形斜边的平方等于两腰的平方和。

可以用符号语言表述为：设直角三角形$ABC$,其中$\angle C=90^\circ$


    \end{thm}
    \begin{equation}
        AB^2+BC^2=AC^2
        \label{eq:gg}
    \end{equation}
    满足式\eqref{eq:gg}的整数称为勾股数。第 1 节所说毕
    达哥拉斯学派得到的三元数组就是勾股数。下表列
    出一些较小的勾股数：
    \begin{table}[H]  
        \begin{tabular}{|rrr|}
            \hline
            直角边$a$ & 直角边$b$ & 直角边$c$\\
            \hline
            3 & 4 & 5 \\
            5 & 12 & 13 \\
            \hline
        \end{tabular}%
        \qquad
        ($a^2+b^2=c^2$)    
    \end{table}
    \bibliography{math}

\end{document}